% from ega4-2-fr.tex
\oldpage[IV-2]{136}

toute composante irréductible de $f^{-1}(V(\mathfrak p))$ domine
$V(\mathfrak p)$, et par suite la condition de (6.1.1) est satisfaite, d'où la
conclusion.

\begin{corollary}[6.1.3]
\EGAReaderTarget{ega4.c04.cor.6-1-3}{IV.2.s6.1.3-fr}
Soient $A$, $B$ deux anneaux locaux noethériens, $\mathfrak m$ l'idéal maximal
de $A$, $k$ son corps résiduel, $\varphi:A\to B$ un homomorphisme local.
Supposons que $\dim(B\otimes_Ak)=0$ (c'est-à-dire (0, 16.2.1) que
$\mathfrak mB$ soit un idéal de définition de $B$). Alors on a
$\dim(B)\leq\dim(A)$. Si en outre il existe un $B$-module de type fini $N$ qui
soit un $A$-module plat et ait un support égal à $\operatorname{Spec}(B)$
(ce qui a lieu en particulier si $B$ est un $A$-module plat), on a
$\dim(B)=\dim(A)$.
\end{corollary}

La première assertion résulte de (5.5.1) et la seconde de (6.1.2).

\begin{corollary}[6.1.4]
\label{IV.2.s6.1.4-fr}
Soient $X$, $Y$ deux préschémas localement noethériens, $f:X\to Y$ un
morphisme surjectif. Pour toute partie fermée $Z$ de $Y$, on a alors
\begin{equation}
\operatorname{codim}(f^{-1}(Z),X)\leq\operatorname{codim}(Z,Y).
\tag{6.1.4.1}\label{IV.2.s6.1.4.1-fr}
\end{equation}
En outre, si $f$ est quasi-plat (2.3.3) les deux membres de (6.1.4.1) sont
égaux.
\end{corollary}

En effet, si $y$ est un point maximal de $Z$ et $x$ un point maximal de la
fibre $f^{-1}(y)$, on a
$\dim(\mathcal O_{X,x})\leq\dim(\mathcal O_{Y,y})$ en vertu de (5.5.2);
l'inégalité (6.1.4.1) résulte alors de (5.1.2.1) et de (5.1.3.1) et du fait
que $f$ est surjectif. Si en outre $f$ est quasi-plat, on sait (2.3.4) que
toute composante irréductible de $f^{-1}(Z)$ domine une composante
irréductible de $Z$; les points maximaux de $f^{-1}(Z)$ sont donc les points
maximaux des fibres $f^{-1}(y)$, où $y$ parcourt l'ensemble des points maximaux
de $Z$ $(0_{\mathrm I}, 2.1.8)$, et en chacun de ces points maximaux $x$ on a
$\dim(\mathcal O_x\otimes_{\mathcal O_y}k(x))=0$; comme en outre
$\mathcal O_x$ est un $\mathcal O_y$-module plat, on a
$\dim(\mathcal O_x)=\dim(\mathcal O_y)$ en vertu de (6.1.1), d'où l'égalité
dans (6.1.4.1), compte tenu de ce que $f$ est surjectif.

La proposition (6.1.1) admet la réciproque partielle suivante:

\begin{proposition}[6.1.5]
\EGAReaderTarget{ega4.c04.prop.6-1-5}{IV.2.s6.1.5-fr}
Soient $A$, $B$ deux anneaux locaux noethériens, $k$ le corps résiduel de
$A$, $\varphi:A\to B$ un homomorphisme local, $M$ un $B$-module de type fini.
On suppose que:
\begin{enumerate}
\item[1\textsuperscript{o}] $A$ est un anneau régulier.
\item[2\textsuperscript{o}] $M$ est un $B$-module de Cohen-Macaulay.
\item[3\textsuperscript{o}] On a
$\dim_B(M)=\dim(A)+\dim_{B\otimes_Ak}(M\otimes_Ak)$.
\end{enumerate}
Alors $M$ est un $A$-module plat.
\end{proposition}

Procédons par récurrence sur $n=\dim(A)$. Si $n=0$, $A$ est un corps
puisqu'il est régulier (0, 17.1.4) et l'assertion est triviale. Supposons
$n>0$; soit $\mathfrak m$ l'idéal maximal de $A$, et soit
$x\in\mathfrak m$ un élément n'appartenant pas à $\mathfrak m^2$; on sait
alors (0, 17.1.8) que $A'=A/xA$ est régulier et
$\dim(A')=\dim(A)-1$. Posons
\[
B'=B/xB,\qquad M'=M/xM=M\otimes_AA';
\]
on a donc
\[
B'\otimes_{A'}k=B\otimes_Ak,\qquad
M'\otimes_{A'}k=M\otimes_Ak
\]
et en vertu de (5.5.1.2) on a
\[
\dim_{B'}(M')\leq\dim(A')+
\dim_{B'\otimes_{A'}k}(M'\otimes_{A'}k);
\]

